\section{Contexto}

\begin{frame}{El problema: intercambiar límite e integral}
    Si $f_n \to f$ en $[a,b]$, ¿cuándo se cumple
    \[
        \lim_{n\to\infty}\int_a^b f_n(x)\,dx = \int_a^b \lim_{n\to\infty} f_n(x)\,dx\,?
    \]
    \begin{itemize}
        \item La \textbf{convergencia uniforme} es suficiente:
        \[
            \abs{\int_a^b f_n-\int_a^b f}\le (b-a)\sup_{x\in[a,b]}\abs{f_n(x)-f(x)}.
        \]
        \item Pero es demasiado exigente: $x^n$ no converge uniformemente en $[0,1]$ y aun así $\int_0^1 x^n\,dx\to 0=\int_0^1 \lim_{n\to\infty} x^n\,dx$.
        \item La \textbf{convergencia puntual} por sí sola no basta.
        \item \alert{Arzelà:} basta la convergencia puntual con una cota uniforme, si el límite es integrable.
    \end{itemize}
\end{frame}


\begin{frame}{Más contexto}
    Tal vez contexto historico, motivación y demás
\end{frame}
